\chapter{Les filtres : une description, douze plateformes}

Un filtre est une condition sous laquelle certains réglages s'appliquent. C'est
le mécanisme qui permet à un seul \texttt{.jenga} de décrire un programme pour
Windows, Linux, Android et le Web sans se dédoubler.

\section{La forme}

\begin{nkcode}{}
with filter("system:Windows"):
    defines(["WIN32_LEAN_AND_MEAN"])
    links(["user32", "gdi32", "opengl32"])

with filter("config:Release"):
    defines(["NDEBUG"])
    optimize("Speed")
    symbols(False)
\end{nkcode}

\begin{nkretenir}
\textbf{Un filtre est un gestionnaire de contexte, comme \texttt{project}.} Tout
ce qui est écrit dedans est \emph{étiqueté} de cette condition et rangé de côté ;
au moment de construire, Jenga ne retient que les étiquettes qui correspondent à
la cible du jour.

Rien n'est évalué à l'écriture : \texttt{with filter("system:Linux")} sur une
machine Windows ne saute pas le bloc --- il l'enregistre. C'est ce qui permet de
décrire douze plateformes depuis une seule.
\end{nkretenir}

\section{Les huit familles de conditions}

Elles sont dans l'évaluateur, et ce sont les seules :

\begin{center}
\begin{tabular}{@{}lll@{}}
\toprule
\textbf{Préfixe} & \textbf{Exemple} & \textbf{Teste} \\
\midrule
\texttt{system:} & \texttt{system:Windows} & le système visé \\
\texttt{config:} & \texttt{config:Debug} & la configuration \\
\texttt{arch:} & \texttt{arch:arm64} & l'architecture \\
\texttt{platform:} & \texttt{platform:Windows-x86\_64} & système + architecture \\
\texttt{action:} & \texttt{action:gen-cmake} & la commande en cours \\
\texttt{options:} & \texttt{options:linux-backend=xlib} & une option déclarée \\
\texttt{kind:} & \texttt{kind:StaticLib} & le genre du projet \\
\texttt{language:} & \texttt{language:C++} & le langage \\
\texttt{toolset:} & \texttt{toolset:clang} & la famille de compilateur \\
\bottomrule
\end{tabular}
\end{center}

\begin{nkretenir}
\texttt{configurations:}, \texttt{configuration:}, \texttt{config:} et
\texttt{cfg:} sont \textbf{quatre orthographes du même test}. De même
\texttt{architecture:} et \texttt{arch:}.

C'est commode et c'est un piège de recherche : chercher \texttt{config:} dans un
dépôt rate les fichiers qui écrivent \texttt{cfg:}. \textbf{Fixez une orthographe
par équipe.}
\end{nkretenir}

\section{Combiner}

\begin{nkcode}{Applications/NkDames/NkDames.jenga -- reel}
with filter("system:Windows && !options:windows-runtime=uwp && !system:XboxSeries"):
    ...
with filter("system:Linux && options:linux-backend=xlib || "
            "system:Linux && !options:linux-backend && !options:headless"):
    ...
\end{nkcode}

\begin{center}
\begin{tabular}{@{}ll@{}}
\toprule
\texttt{\&\&} & les deux conditions \\
\texttt{||} & l'une ou l'autre \\
\texttt{!} & la négation \\
\bottomrule
\end{tabular}
\end{center}

\begin{nkretenir}
\textbf{La seconde expression mérite d'être lue lentement} : elle dit « backend
xlib demandé explicitement, \emph{ou} aucun backend demandé et pas de mode sans
écran ».

C'est ainsi qu'on écrit un \textbf{défaut} : le second membre attrape le cas où
l'utilisateur n'a rien choisi. Sans lui, ne rien demander ne donnerait
\emph{aucun} backend --- et l'échec sortirait à l'édition de liens.
\end{nkretenir}

\begin{nkretenir}
Une liste ou un ensemble passés à \texttt{filter} sont joints par \texttt{\&\&} :

\begin{nkcode}{Jenga/Core/Api.py:597-609}
if isinstance(expr, (list, tuple)):
    parts = [str(part).strip() for part in expr if str(part).strip()]
    return " && ".join(parts)
\end{nkcode}

\texttt{filter(["system:Windows", "config:Release"])} vaut donc
\texttt{"system:Windows \&\& config:Release"}. \alerte{} Un \emph{ensemble} est en plus
\textbf{trié} --- pour que la clé soit stable, l'ordre d'un ensemble Python ne
l'étant pas.
\end{nkretenir}

\section{Les options : vos propres conditions}

\begin{nkcode}{}
newoption({
    "trigger": "linux-backend",
    "value": "xlib",
    "description": "Backend fenetrage sous Linux",
    "allowed": ["xlib", "xcb", "wayland", "headless"],
})
\end{nkcode}

\begin{nkcode}{Puis, en ligne de commande}
jenga build --platform linux --linux-backend=wayland
\end{nkcode}

\begin{nkretenir}
\textbf{\texttt{allowed} n'est pas décoratif} : il fait refuser une valeur
inconnue \emph{au chargement}, avec un message qui énumère les valeurs admises
--- exactement comme l'erreur de \texttt{-{}-kind} du chapitre 3.

Sans cette liste, une faute de frappe donne une option silencieusement ignorée,
et vous construisez pour un backend que vous n'avez pas demandé.
\end{nkretenir}

\section{Le piège du filtre qui n'attrape rien}

\begin{nkpiege}
\textbf{Un filtre qui ne correspond à rien ne produit AUCUN message.} C'est
normal --- c'est même son travail --- mais cela signifie qu'une faute de frappe
dans la condition est indiscernable d'une condition légitimement non remplie.

\texttt{with filter("system:windows")} en minuscules fonctionne, l'évaluateur
compare en minuscules. Mais \texttt{with filter("systeme:Windows")} --- en
français --- ne correspond à \emph{aucune} famille connue, ne déclenche rien, et
ne se plaint pas.

Le symptôme sera une bibliothèque manquante à l'édition de liens.
\end{nkpiege}

\begin{nkretenir}[Comment on l'attrape]
\textbf{\texttt{jenga compile-flags} affiche les drapeaux réellement retenus}
pour chaque projet. Si votre \texttt{defines} n'y est pas, son filtre n'a pas
mordu.

C'est le seul contrôle qui distingue « la condition est fausse aujourd'hui » de
« la condition ne peut jamais être vraie » --- en la testant sur une cible où
elle \emph{devrait} être vraie.
\end{nkretenir}

\section{Ce qu'un filtre ne peut pas faire}

\begin{nkretenir}
\textbf{Un filtre ne crée pas de projet.} Il modifie les réglages d'un projet
existant. Pour ne déclarer un projet que sur certaines plateformes, employez du
\textbf{Python} --- un \texttt{if} autour du \texttt{with project} --- puisque le
fichier est exécuté.

Les deux mécanismes ont l'air interchangeables et ne le sont pas : le
\texttt{if} décide \emph{au chargement}, sur la machine courante ; le filtre
décide \emph{à la construction}, sur la cible. Un \texttt{if} sur
\texttt{sys.platform} rendrait votre descripteur incapable de décrire une
construction croisée.
\end{nkretenir}

\begin{nkbilan}
Un filtre étiquette des réglages d'une condition ; rien n'est évalué à
l'écriture, ce qui permet de décrire douze plateformes depuis une seule. Neuf
familles de conditions existent, certaines avec plusieurs orthographes --- fixez
la vôtre. Les expressions se combinent, et un second membre sert souvent à
attraper le cas « rien n'a été demandé ». Vos propres options se déclarent par
\texttt{newoption}, dont la liste \texttt{allowed} transforme une faute de frappe
en erreur au lieu d'un silence. Enfin un filtre ne crée pas de projet : cela,
c'est le rôle du Python --- mais un \texttt{if} décide sur la machine courante là
où un filtre décide sur la cible.
\end{nkbilan}

\begin{nkquestions}
\begin{enumerate}
\item Que fait Jenga d'un bloc \texttt{filter} dont la condition est fausse sur
      la machine courante ?
\item Citez cinq familles de conditions.
\item Que vaut \texttt{filter(["a", "b"])} ?
\item Pourquoi une expression de filtre comporte-t-elle souvent un second membre
      avec des négations ?
\item Quand faut-il un \texttt{if} Python plutôt qu'un filtre ?
\end{enumerate}
\end{nkquestions}

\begin{nkpratique}
Reprenez votre projet et donnez-lui trois comportements :

\begin{enumerate}
\item en \texttt{Debug}, une définition \texttt{MODE\_BAVARD} ;
\item en \texttt{Release}, l'optimisation en vitesse et pas de symboles ;
\item une option à vous, \texttt{-{}-couleur=} avec trois valeurs admises, qui
      pose une définition différente selon la valeur.
\end{enumerate}

Vérifiez les trois par \texttt{jenga compile-flags}, pas par la construction :
c'est plus rapide et plus précis.
\end{nkpratique}

\begin{nkdemo}
Prouvez qu'un filtre mal orthographié se tait.

Écrivez deux filtres qui posent chacun une définition : l'un correct
(\texttt{system:Windows}), l'autre volontairement faux
(\texttt{systeme:Windows}). Construisez sur Windows.

\textbf{Ce que la démonstration doit établir}, en trois relevés :
\begin{enumerate}
\item la construction \textbf{réussit} ;
\item aucun message ne mentionne le filtre fautif ;
\item \texttt{jenga compile-flags} montre la première définition et pas la
      seconde.
\end{enumerate}

\alerte{} Les deux premiers points sont le cœur de la démonstration. Si vous ne relevez
que le troisième, vous prouvez que l'outil sait montrer les drapeaux --- pas que
le défaut est silencieux.
\end{nkdemo}

\begin{nklogique}
Un projet lie \texttt{X11} sur Linux. Sur la machine d'un collègue, l'édition de
liens échoue en disant que \texttt{X11} est introuvable.

\begin{enumerate}
\item Énumérez trois causes, en distinguant celles qui vivent dans le
      \texttt{.jenga} de celles qui vivent sur la machine.
\item Pour chacune, dites ce que montrerait \texttt{jenga compile-flags} --- et
      pour laquelle il ne montrerait \emph{rien de particulier}.
\item Cette dernière est la plus fréquente. Expliquez pourquoi l'outil ne peut
      pas la voir, et ce qui la verrait.
\end{enumerate}
\end{nklogique}
